%%%%%%%%%%%%%%%%%%%%%%%%%%%%%%%%%%%%%%%%%%%%%%%%%%%%%%%%%%%%
% 9.6 INTEGRAL EQUATIONS
% The Composition of Advanced Mathematics
%%%%%%%%%%%%%%%%%%%%%%%%%%%%%%%%%%%%%%%%%%%%%%%%%%%%%%%%%%%%

\section{Integral Equations}
\label{sec:integral-equations}

\prerequisite{
Integration, improper integrals, ordinary differential equations,
boundary-value problems, linear algebra, eigenvalues and eigenvectors,
transform methods, convolution, Green's functions, sequences and series,
and basic functional analysis concepts.
}

\learningobjectives{
By the end of this section, the reader should be able to:
\begin{itemize}
    \item define an integral equation;
    \item identify the unknown function and kernel;
    \item distinguish Fredholm and Volterra integral equations;
    \item distinguish integral equations of the first and second kind;
    \item distinguish homogeneous and nonhomogeneous integral equations;
    \item distinguish linear and nonlinear integral equations;
    \item recognize convolution integral equations;
    \item understand integral operators;
    \item rewrite integral equations using operator notation;
    \item solve elementary equations with separable kernels;
    \item solve equations with degenerate kernels;
    \item understand iterative substitution;
    \item derive successive approximations;
    \item understand Neumann series;
    \item define iterated kernels;
    \item understand resolvent kernels;
    \item solve Volterra equations by differentiation when appropriate;
    \item solve convolution-type equations using Laplace transforms;
    \item understand the relation between integral equations and ODEs;
    \item understand the relation between integral equations and
          boundary-value problems;
    \item understand Green's-function formulations;
    \item recognize eigenvalue integral equations;
    \item understand symmetric kernels;
    \item recognize the Fredholm alternative conceptually;
    \item understand compact integral operators conceptually;
    \item recognize nonlinear integral equations;
    \item understand Hammerstein and Urysohn equations conceptually;
    \item recognize singular integral equations;
    \item understand Abel-type integral equations conceptually;
    \item approximate integral equations numerically;
    \item understand quadrature-based methods;
    \item understand the Nystr\"om method;
    \item understand collocation and Galerkin methods conceptually;
    \item connect integral equations with ODEs, PDEs, transforms,
          functional analysis, and numerical mathematics.
\end{itemize}
}

%%%%%%%%%%%%%%%%%%%%%%%%%%%%%%%%%%%%%%%%%%%%%%%%%%%%%%%%%%%%
\subsection{What Is an Integral Equation?}
\label{subsec:what-is-integral-equation}
%%%%%%%%%%%%%%%%%%%%%%%%%%%%%%%%%%%%%%%%%%%%%%%%%%%%%%%%%%%%

An integral equation is an equation in which the unknown function appears
inside an integral.

A simple example is

\[
\boxed{
u(x)
=
1
+
\int_0^x
u(t)\,dt.
}
\]

The unknown function is

\[
u.
\]

Unlike an ordinary algebraic equation, the value \(u(x)\) depends on an
integral involving the entire function over some interval.

%%%%%%%%%%%%%%%%%%%%%%%%%%%%%%%%%%%%%%%%%%%%%%%%%%%%%%%%%%%%
\subsection{General Linear Integral Equation}
\label{subsec:general-linear-integral-equation}
%%%%%%%%%%%%%%%%%%%%%%%%%%%%%%%%%%%%%%%%%%%%%%%%%%%%%%%%%%%%

A common linear integral equation has form

\[
\boxed{
u(x)
=
f(x)
+
\lambda
\int_a^b
K(x,t)u(t)\,dt.
}
\]

Here:

\[
u(x)
=
\text{unknown function},
\]

\[
f(x)
=
\text{known forcing function},
\]

\[
K(x,t)
=
\text{kernel},
\]

and

\[
\lambda
=
\text{parameter}.
\]

%%%%%%%%%%%%%%%%%%%%%%%%%%%%%%%%%%%%%%%%%%%%%%%%%%%%%%%%%%%%
\subsection{Kernel}
\label{subsec:integral-equation-kernel}
%%%%%%%%%%%%%%%%%%%%%%%%%%%%%%%%%%%%%%%%%%%%%%%%%%%%%%%%%%%%

\begin{definitionbox}{Kernel}
The kernel of a linear integral equation is the known function \(K(x,t)\)
that determines how values of the unknown function \(u(t)\) influence
the equation at the point \(x\).
\end{definitionbox}

The kernel can encode:

\[
\boxed{
\text{interaction},
}
\]

\[
\boxed{
\text{memory},
}
\]

\[
\boxed{
\text{spatial influence},
}
\]

or

\[
\boxed{
\text{propagation}.
}
\]

%%%%%%%%%%%%%%%%%%%%%%%%%%%%%%%%%%%%%%%%%%%%%%%%%%%%%%%%%%%%
\subsection{Fredholm Integral Equations}
\label{subsec:fredholm-integral-equations}
%%%%%%%%%%%%%%%%%%%%%%%%%%%%%%%%%%%%%%%%%%%%%%%%%%%%%%%%%%%%

A Fredholm integral equation has fixed integration limits.

A Fredholm equation of the second kind is

\[
\boxed{
u(x)
=
f(x)
+
\lambda
\int_a^b
K(x,t)u(t)\,dt.
}
\]

The interval

\[
[a,b]
\]

does not depend on \(x\).

%%%%%%%%%%%%%%%%%%%%%%%%%%%%%%%%%%%%%%%%%%%%%%%%%%%%%%%%%%%%
\subsection{Volterra Integral Equations}
\label{subsec:volterra-integral-equations}
%%%%%%%%%%%%%%%%%%%%%%%%%%%%%%%%%%%%%%%%%%%%%%%%%%%%%%%%%%%%

A Volterra integral equation has a variable integration limit.

A Volterra equation of the second kind is

\[
\boxed{
u(x)
=
f(x)
+
\lambda
\int_a^x
K(x,t)u(t)\,dt.
}
\]

The integral depends only on earlier values

\[
t\leq x.
\]

Thus Volterra equations naturally represent causal or memory processes.

%%%%%%%%%%%%%%%%%%%%%%%%%%%%%%%%%%%%%%%%%%%%%%%%%%%%%%%%%%%%
\subsection{Fredholm Versus Volterra}
\label{subsec:fredholm-versus-volterra}
%%%%%%%%%%%%%%%%%%%%%%%%%%%%%%%%%%%%%%%%%%%%%%%%%%%%%%%%%%%%

The essential distinction is:

\[
\boxed{
\text{Fredholm}
\rightarrow
\text{fixed integration interval},
}
\]

while

\[
\boxed{
\text{Volterra}
\rightarrow
\text{variable integration limit}.
}
\]

Volterra equations are often more directly connected with initial-value
problems.

%%%%%%%%%%%%%%%%%%%%%%%%%%%%%%%%%%%%%%%%%%%%%%%%%%%%%%%%%%%%
\subsection{Integral Equations of the First Kind}
\label{subsec:first-kind-integral-equations}
%%%%%%%%%%%%%%%%%%%%%%%%%%%%%%%%%%%%%%%%%%%%%%%%%%%%%%%%%%%%

A first-kind linear integral equation has form

\[
\boxed{
f(x)
=
\int_a^b
K(x,t)u(t)\,dt.
}
\]

The unknown function appears only under the integral.

First-kind equations can be mathematically more difficult and numerically
ill-conditioned.

%%%%%%%%%%%%%%%%%%%%%%%%%%%%%%%%%%%%%%%%%%%%%%%%%%%%%%%%%%%%
\subsection{Integral Equations of the Second Kind}
\label{subsec:second-kind-integral-equations}
%%%%%%%%%%%%%%%%%%%%%%%%%%%%%%%%%%%%%%%%%%%%%%%%%%%%%%%%%%%%

A second-kind equation has the unknown both inside and outside the
integral:

\[
\boxed{
u(x)
=
f(x)
+
\lambda
\int_a^b
K(x,t)u(t)\,dt.
}
\]

Second-kind equations often have better stability properties than
first-kind equations.

%%%%%%%%%%%%%%%%%%%%%%%%%%%%%%%%%%%%%%%%%%%%%%%%%%%%%%%%%%%%
\subsection{Third-Kind Integral Equations}
\label{subsec:third-kind-integral-equations}
%%%%%%%%%%%%%%%%%%%%%%%%%%%%%%%%%%%%%%%%%%%%%%%%%%%%%%%%%%%%

A more general form is

\[
\boxed{
a(x)u(x)
=
f(x)
+
\lambda
\int_a^b
K(x,t)u(t)\,dt.
}
\]

If

\[
a(x)
\]

can vanish, the equation may display behavior intermediate between first-
and second-kind equations.

%%%%%%%%%%%%%%%%%%%%%%%%%%%%%%%%%%%%%%%%%%%%%%%%%%%%%%%%%%%%
\subsection{Homogeneous Integral Equations}
\label{subsec:homogeneous-integral-equations}
%%%%%%%%%%%%%%%%%%%%%%%%%%%%%%%%%%%%%%%%%%%%%%%%%%%%%%%%%%%%

A homogeneous second-kind equation has

\[
f(x)=0.
\]

Thus,

\[
\boxed{
u(x)
=
\lambda
\int_a^b
K(x,t)u(t)\,dt.
}
\]

Nonzero solutions exist only for special parameter values in many
important problems.

%%%%%%%%%%%%%%%%%%%%%%%%%%%%%%%%%%%%%%%%%%%%%%%%%%%%%%%%%%%%
\subsection{Nonhomogeneous Integral Equations}
\label{subsec:nonhomogeneous-integral-equations}
%%%%%%%%%%%%%%%%%%%%%%%%%%%%%%%%%%%%%%%%%%%%%%%%%%%%%%%%%%%%

If

\[
f(x)\neq0,
\]

then

\[
\boxed{
u(x)
=
f(x)
+
\lambda
\int_a^b
K(x,t)u(t)\,dt
}
\]

is nonhomogeneous.

This is analogous to a forced differential equation.

%%%%%%%%%%%%%%%%%%%%%%%%%%%%%%%%%%%%%%%%%%%%%%%%%%%%%%%%%%%%
\subsection{Linear Versus Nonlinear Integral Equations}
\label{subsec:linear-versus-nonlinear-integral-equations}
%%%%%%%%%%%%%%%%%%%%%%%%%%%%%%%%%%%%%%%%%%%%%%%%%%%%%%%%%%%%

A linear equation contains the unknown linearly:

\[
\boxed{
u(x)
=
f(x)
+
\int
K(x,t)u(t)\,dt.
}
\]

A nonlinear example is

\[
\boxed{
u(x)
=
f(x)
+
\int_a^b
K(x,t)
[u(t)]^2
\,dt.
}
\]

Another is

\[
\boxed{
u(x)
=
f(x)
+
\int_a^b
K(x,t)
\sin(u(t))
\,dt.
}
\]

%%%%%%%%%%%%%%%%%%%%%%%%%%%%%%%%%%%%%%%%%%%%%%%%%%%%%%%%%%%%
\subsection{Integral Operators}
\label{subsec:integral-operators}
%%%%%%%%%%%%%%%%%%%%%%%%%%%%%%%%%%%%%%%%%%%%%%%%%%%%%%%%%%%%

Define an integral operator

\[
\boxed{
(Tu)(x)
=
\int_a^b
K(x,t)u(t)\,dt.
}
\]

Then a second-kind equation becomes

\[
\boxed{
u
=
f+\lambda Tu.
}
\]

Equivalently,

\[
\boxed{
(I-\lambda T)u
=
f.
}
\]

This operator form connects integral equations with linear algebra.

%%%%%%%%%%%%%%%%%%%%%%%%%%%%%%%%%%%%%%%%%%%%%%%%%%%%%%%%%%%%
\subsection{Analogy with Linear Algebra}
\label{subsec:integral-equation-linear-algebra-analogy}
%%%%%%%%%%%%%%%%%%%%%%%%%%%%%%%%%%%%%%%%%%%%%%%%%%%%%%%%%%%%

Compare

\[
\boxed{
(I-\lambda T)u=f
}
\]

with the matrix equation

\[
\boxed{
(I-\lambda A)\mathbf{x}
=
\mathbf{b}.
}
\]

The integral operator \(T\) plays a role analogous to a matrix.

The unknown function \(u\) plays a role analogous to a vector.

This analogy becomes fundamental in functional analysis.

%%%%%%%%%%%%%%%%%%%%%%%%%%%%%%%%%%%%%%%%%%%%%%%%%%%%%%%%%%%%
\subsection{Separable Kernels}
\label{subsec:separable-kernels}
%%%%%%%%%%%%%%%%%%%%%%%%%%%%%%%%%%%%%%%%%%%%%%%%%%%%%%%%%%%%

A kernel is separable if it can be written

\[
\boxed{
K(x,t)
=
g(x)h(t).
}
\]

More generally, a finite-rank or degenerate kernel has form

\[
\boxed{
K(x,t)
=
\sum_{j=1}^{m}
g_j(x)h_j(t).
}
\]

Such equations can often be reduced to finite systems of algebraic
equations.

%%%%%%%%%%%%%%%%%%%%%%%%%%%%%%%%%%%%%%%%%%%%%%%%%%%%%%%%%%%%
\subsection{Worked Example: Separable Kernel}
\label{subsec:worked-separable-kernel}
%%%%%%%%%%%%%%%%%%%%%%%%%%%%%%%%%%%%%%%%%%%%%%%%%%%%%%%%%%%%

\begin{workedexamplebox}{Reducing an Integral Equation to Algebra}

Solve

\[
u(x)
=
1
+
x
\int_0^1
t\,u(t)\,dt.
\]

Define

\[
C
=
\int_0^1
t\,u(t)\,dt.
\]

Then

\[
u(x)=1+xC.
\]

Substitute into the definition of \(C\):

\[
C
=
\int_0^1
t(1+tC)\,dt.
\]

Thus,

\[
C
=
\int_0^1t\,dt
+
C
\int_0^1t^2\,dt.
\]

Therefore,

\[
C
=
\frac12
+
\frac13C.
\]

Hence,

\[
\frac23C
=
\frac12,
\]

so

\[
C
=
\frac34.
\]

Therefore,

\begin{answerbox}
\[
\boxed{
u(x)
=
1+\frac34x.
}
\]
\end{answerbox}

\end{workedexamplebox}

%%%%%%%%%%%%%%%%%%%%%%%%%%%%%%%%%%%%%%%%%%%%%%%%%%%%%%%%%%%%
\subsection{Degenerate Kernels}
\label{subsec:degenerate-kernels}
%%%%%%%%%%%%%%%%%%%%%%%%%%%%%%%%%%%%%%%%%%%%%%%%%%%%%%%%%%%%

Suppose

\[
K(x,t)
=
\sum_{j=1}^{m}
g_j(x)h_j(t).
\]

Then

\[
u(x)
=
f(x)
+
\lambda
\sum_{j=1}^{m}
g_j(x)
\int_a^b
h_j(t)u(t)\,dt.
\]

Define unknown constants

\[
\boxed{
c_j
=
\int_a^b
h_j(t)u(t)\,dt.
}
\]

Then

\[
\boxed{
u(x)
=
f(x)
+
\lambda
\sum_{j=1}^{m}
c_jg_j(x).
}
\]

The \(c_j\) satisfy a finite linear system.

%%%%%%%%%%%%%%%%%%%%%%%%%%%%%%%%%%%%%%%%%%%%%%%%%%%%%%%%%%%%
\subsection{Successive Approximation}
\label{subsec:successive-approximation-integral-equations}
%%%%%%%%%%%%%%%%%%%%%%%%%%%%%%%%%%%%%%%%%%%%%%%%%%%%%%%%%%%%

Consider

\[
u=f+\lambda Tu.
\]

Begin with

\[
\boxed{
u_0=f.
}
\]

Then define iteratively

\[
\boxed{
u_{n+1}
=
f+\lambda Tu_n.
}
\]

This is called successive approximation or Picard iteration.

%%%%%%%%%%%%%%%%%%%%%%%%%%%%%%%%%%%%%%%%%%%%%%%%%%%%%%%%%%%%
\subsection{First Few Successive Approximations}
\label{subsec:first-successive-approximations}
%%%%%%%%%%%%%%%%%%%%%%%%%%%%%%%%%%%%%%%%%%%%%%%%%%%%%%%%%%%%

Starting with

\[
u_0=f,
\]

we obtain

\[
u_1
=
f+\lambda Tf,
\]

\[
u_2
=
f+\lambda Tf+\lambda^2T^2f,
\]

and

\[
u_3
=
f+\lambda Tf+\lambda^2T^2f+\lambda^3T^3f.
\]

Thus the pattern suggests

\[
\boxed{
u
=
f
+
\lambda Tf
+
\lambda^2T^2f
+
\lambda^3T^3f
+
\cdots.
}
\]

%%%%%%%%%%%%%%%%%%%%%%%%%%%%%%%%%%%%%%%%%%%%%%%%%%%%%%%%%%%%
\subsection{Neumann Series}
\label{subsec:neumann-series-integral-equations}
%%%%%%%%%%%%%%%%%%%%%%%%%%%%%%%%%%%%%%%%%%%%%%%%%%%%%%%%%%%%

The formal operator series

\[
\boxed{
(I-\lambda T)^{-1}
=
I+\lambda T+\lambda^2T^2+\cdots
}
\]

is called a Neumann series.

Thus,

\[
\boxed{
u
=
\sum_{n=0}^{\infty}
\lambda^nT^nf.
}
\]

This is analogous to the matrix geometric series

\[
(I-A)^{-1}
=
I+A+A^2+\cdots.
\]

%%%%%%%%%%%%%%%%%%%%%%%%%%%%%%%%%%%%%%%%%%%%%%%%%%%%%%%%%%%%
\subsection{Convergence of the Neumann Series}
\label{subsec:neumann-series-convergence}
%%%%%%%%%%%%%%%%%%%%%%%%%%%%%%%%%%%%%%%%%%%%%%%%%%%%%%%%%%%%

A sufficient operator-norm condition is

\[
\boxed{
|\lambda|
\|T\|
<
1.
}
\]

Then

\[
\boxed{
(I-\lambda T)^{-1}
}
\]

exists and the Neumann series converges in the relevant operator norm.

The condition is sufficient, not always necessary.

%%%%%%%%%%%%%%%%%%%%%%%%%%%%%%%%%%%%%%%%%%%%%%%%%%%%%%%%%%%%
\subsection{Iterated Kernels}
\label{subsec:iterated-kernels}
%%%%%%%%%%%%%%%%%%%%%%%%%%%%%%%%%%%%%%%%%%%%%%%%%%%%%%%%%%%%

Define the first kernel

\[
\boxed{
K_1(x,t)=K(x,t).
}
\]

The second iterated kernel is

\[
\boxed{
K_2(x,t)
=
\int_a^b
K(x,s)K(s,t)\,ds.
}
\]

Recursively,

\[
\boxed{
K_{n+1}(x,t)
=
\int_a^b
K(x,s)K_n(s,t)\,ds.
}
\]

%%%%%%%%%%%%%%%%%%%%%%%%%%%%%%%%%%%%%%%%%%%%%%%%%%%%%%%%%%%%
\subsection{Neumann Series in Kernel Form}
\label{subsec:neumann-kernel-form}
%%%%%%%%%%%%%%%%%%%%%%%%%%%%%%%%%%%%%%%%%%%%%%%%%%%%%%%%%%%%

The solution can formally be written

\[
u(x)
=
f(x)
+
\sum_{n=1}^{\infty}
\lambda^n
\int_a^b
K_n(x,t)f(t)\,dt.
\]

This motivates the resolvent kernel.

%%%%%%%%%%%%%%%%%%%%%%%%%%%%%%%%%%%%%%%%%%%%%%%%%%%%%%%%%%%%
\subsection{Resolvent Kernel}
\label{subsec:resolvent-kernel}
%%%%%%%%%%%%%%%%%%%%%%%%%%%%%%%%%%%%%%%%%%%%%%%%%%%%%%%%%%%%

Define

\[
\boxed{
R(x,t;\lambda)
=
K_1(x,t)
+
\lambda K_2(x,t)
+
\lambda^2K_3(x,t)
+
\cdots.
}
\]

When the series converges,

\[
\boxed{
u(x)
=
f(x)
+
\lambda
\int_a^b
R(x,t;\lambda)
f(t)\,dt.
}
\]

The resolvent kernel plays a role analogous to an inverse matrix.

%%%%%%%%%%%%%%%%%%%%%%%%%%%%%%%%%%%%%%%%%%%%%%%%%%%%%%%%%%%%
\subsection{Volterra Equations and Successive Approximation}
\label{subsec:volterra-successive-approximation}
%%%%%%%%%%%%%%%%%%%%%%%%%%%%%%%%%%%%%%%%%%%%%%%%%%%%%%%%%%%%

For

\[
u(x)
=
f(x)
+
\lambda
\int_a^x
K(x,t)u(t)\,dt,
\]

successive approximation is particularly effective.

Because the integration region is triangular,

\[
a\leq t\leq x,
\]

iterated integrals often introduce factorial denominators that improve
convergence.

%%%%%%%%%%%%%%%%%%%%%%%%%%%%%%%%%%%%%%%%%%%%%%%%%%%%%%%%%%%%
\subsection{Simple Volterra Equation}
\label{subsec:simple-volterra-equation}
%%%%%%%%%%%%%%%%%%%%%%%%%%%%%%%%%%%%%%%%%%%%%%%%%%%%%%%%%%%%

Consider

\[
\boxed{
u(x)
=
1
+
\int_0^x
u(t)\,dt.
}
\]

Differentiate both sides:

\[
u'(x)
=
u(x).
\]

From the original equation,

\[
u(0)=1.
\]

Thus the equivalent IVP is

\[
\boxed{
u'=u,
\qquad
u(0)=1.
}
\]

Therefore,

\[
\boxed{
u(x)=e^x.
}
\]

%%%%%%%%%%%%%%%%%%%%%%%%%%%%%%%%%%%%%%%%%%%%%%%%%%%%%%%%%%%%
\subsection{Worked Example: Volterra Equation by Differentiation}
\label{subsec:worked-volterra-differentiation}
%%%%%%%%%%%%%%%%%%%%%%%%%%%%%%%%%%%%%%%%%%%%%%%%%%%%%%%%%%%%

\begin{workedexamplebox}{Converting an Integral Equation to an ODE}

Solve

\[
u(x)
=
x
+
\int_0^x
u(t)\,dt.
\]

Differentiate:

\[
u'(x)
=
1+u(x).
\]

From the original equation,

\[
u(0)=0.
\]

Thus,

\[
u'-u=1.
\]

The general solution is

\[
u(x)
=
Ce^x-1.
\]

Using

\[
u(0)=0,
\]

we obtain

\[
C=1.
\]

Therefore,

\begin{answerbox}
\[
\boxed{
u(x)
=
e^x-1.
}
\]
\end{answerbox}

\end{workedexamplebox}

%%%%%%%%%%%%%%%%%%%%%%%%%%%%%%%%%%%%%%%%%%%%%%%%%%%%%%%%%%%%
\subsection{Leibniz Rule for Variable Limits}
\label{subsec:leibniz-rule-integral-equations}
%%%%%%%%%%%%%%%%%%%%%%%%%%%%%%%%%%%%%%%%%%%%%%%%%%%%%%%%%%%%

If

\[
F(x)
=
\int_a^x
K(x,t)u(t)\,dt,
\]

then under suitable regularity,

\[
\boxed{
F'(x)
=
K(x,x)u(x)
+
\int_a^x
K_x(x,t)u(t)\,dt.
}
\]

This formula is useful when converting Volterra equations into
differential equations.

%%%%%%%%%%%%%%%%%%%%%%%%%%%%%%%%%%%%%%%%%%%%%%%%%%%%%%%%%%%%
\subsection{Integral Form of an Initial-Value Problem}
\label{subsec:ivp-integral-form}
%%%%%%%%%%%%%%%%%%%%%%%%%%%%%%%%%%%%%%%%%%%%%%%%%%%%%%%%%%%%

Consider

\[
y'=f(x,y),
\qquad
y(x_0)=y_0.
\]

Integrating gives

\[
\boxed{
y(x)
=
y_0
+
\int_{x_0}^{x}
f(t,y(t))\,dt.
}
\]

Thus every sufficiently regular IVP can be rewritten as a Volterra
integral equation.

%%%%%%%%%%%%%%%%%%%%%%%%%%%%%%%%%%%%%%%%%%%%%%%%%%%%%%%%%%%%
\subsection{Picard Iteration Revisited}
\label{subsec:picard-integral-equation}
%%%%%%%%%%%%%%%%%%%%%%%%%%%%%%%%%%%%%%%%%%%%%%%%%%%%%%%%%%%%

The integral formulation

\[
y(x)
=
y_0
+
\int_{x_0}^{x}
f(t,y(t))\,dt
\]

suggests

\[
\boxed{
y_{n+1}(x)
=
y_0
+
\int_{x_0}^{x}
f(t,y_n(t))\,dt.
}
\]

This is precisely the Picard iteration used in ODE existence theory.

Thus integral equations provide the foundation for an important ODE
existence method.

%%%%%%%%%%%%%%%%%%%%%%%%%%%%%%%%%%%%%%%%%%%%%%%%%%%%%%%%%%%%
\subsection{Convolution Integral Equations}
\label{subsec:convolution-integral-equations}
%%%%%%%%%%%%%%%%%%%%%%%%%%%%%%%%%%%%%%%%%%%%%%%%%%%%%%%%%%%%

A Volterra convolution equation has form

\[
\boxed{
u(t)
=
f(t)
+
\int_0^t
k(t-\tau)u(\tau)\,d\tau.
}
\]

Using convolution notation,

\[
\boxed{
u
=
f+k*u.
}
\]

Laplace transforms convert convolution into multiplication.

%%%%%%%%%%%%%%%%%%%%%%%%%%%%%%%%%%%%%%%%%%%%%%%%%%%%%%%%%%%%
\subsection{Laplace Transform Solution}
\label{subsec:laplace-integral-equation-solution}
%%%%%%%%%%%%%%%%%%%%%%%%%%%%%%%%%%%%%%%%%%%%%%%%%%%%%%%%%%%%

Let

\[
U(s)
=
\mathcal{L}\{u\},
\]

\[
F(s)
=
\mathcal{L}\{f\},
\]

and

\[
K(s)
=
\mathcal{L}\{k\}.
\]

Then

\[
U
=
F+KU.
\]

Therefore,

\[
\boxed{
U(s)
=
\frac{
F(s)
}{
1-K(s)
}.
}
\]

More generally, if a parameter \(\lambda\) appears,

\[
\boxed{
U(s)
=
\frac{
F(s)
}{
1-\lambda K(s)
}.
}
\]

%%%%%%%%%%%%%%%%%%%%%%%%%%%%%%%%%%%%%%%%%%%%%%%%%%%%%%%%%%%%
\subsection{Worked Example: Convolution Equation}
\label{subsec:worked-convolution-integral-equation}
%%%%%%%%%%%%%%%%%%%%%%%%%%%%%%%%%%%%%%%%%%%%%%%%%%%%%%%%%%%%

\begin{workedexamplebox}{Laplace Transform Method}

Solve

\[
u(t)
=
1
+
\int_0^t
u(\tau)\,d\tau.
\]

Here,

\[
k(t)=1.
\]

Thus,

\[
K(s)=\frac1s,
\qquad
F(s)=\frac1s.
\]

The transformed equation is

\[
U(s)
=
\frac1s
+
\frac1sU(s).
\]

Therefore,

\[
U(s)
\left(
1-\frac1s
\right)
=
\frac1s.
\]

Hence,

\[
U(s)
=
\frac1{s-1}.
\]

Thus,

\begin{answerbox}
\[
\boxed{
u(t)=e^t.
}
\]
\end{answerbox}

\end{workedexamplebox}

%%%%%%%%%%%%%%%%%%%%%%%%%%%%%%%%%%%%%%%%%%%%%%%%%%%%%%%%%%%%
\subsection{Boundary-Value Problems and Integral Equations}
\label{subsec:bvp-integral-equations}
%%%%%%%%%%%%%%%%%%%%%%%%%%%%%%%%%%%%%%%%%%%%%%%%%%%%%%%%%%%%

Linear boundary-value problems can often be rewritten as integral
equations using Green's functions.

Consider

\[
\boxed{
L[u]=f
}
\]

with specified boundary conditions.

If \(G(x,\xi)\) is the Green's function for \(L\), then the solution may
have form

\[
\boxed{
u(x)
=
\int_a^b
G(x,\xi)f(\xi)\,d\xi.
}
\]

%%%%%%%%%%%%%%%%%%%%%%%%%%%%%%%%%%%%%%%%%%%%%%%%%%%%%%%%%%%%
\subsection{Green's Function as an Integral Kernel}
\label{subsec:greens-function-integral-kernel}
%%%%%%%%%%%%%%%%%%%%%%%%%%%%%%%%%%%%%%%%%%%%%%%%%%%%%%%%%%%%

A Green's function is an integral-equation kernel adapted to a
differential operator and its boundary conditions.

Symbolically,

\[
\boxed{
L_xG(x,\xi)
=
\delta(x-\xi).
}
\]

The response to distributed forcing is obtained by superposition:

\[
\boxed{
u(x)
=
\int
G(x,\xi)f(\xi)\,d\xi.
}
\]

%%%%%%%%%%%%%%%%%%%%%%%%%%%%%%%%%%%%%%%%%%%%%%%%%%%%%%%%%%%%
\subsection{Integral Form of a Nonlinear Boundary Problem}
\label{subsec:nonlinear-bvp-integral-form}
%%%%%%%%%%%%%%%%%%%%%%%%%%%%%%%%%%%%%%%%%%%%%%%%%%%%%%%%%%%%

Suppose

\[
L[u]
=
N(u),
\]

where \(L\) is linear and \(N\) is nonlinear.

Using a Green's function can produce

\[
\boxed{
u(x)
=
u_0(x)
+
\int_a^b
G(x,t)
N(u(t))
\,dt.
}
\]

Thus nonlinear differential equations may become nonlinear integral
equations.

%%%%%%%%%%%%%%%%%%%%%%%%%%%%%%%%%%%%%%%%%%%%%%%%%%%%%%%%%%%%
\subsection{Eigenvalue Integral Equations}
\label{subsec:eigenvalue-integral-equations}
%%%%%%%%%%%%%%%%%%%%%%%%%%%%%%%%%%%%%%%%%%%%%%%%%%%%%%%%%%%%

Consider

\[
\boxed{
u(x)
=
\lambda
\int_a^b
K(x,t)u(t)\,dt.
}
\]

In operator notation,

\[
\boxed{
u=\lambda Tu.
}
\]

Equivalently,

\[
\boxed{
Tu
=
\frac1\lambda u.
}
\]

Thus \(u\) is an eigenfunction of \(T\).

%%%%%%%%%%%%%%%%%%%%%%%%%%%%%%%%%%%%%%%%%%%%%%%%%%%%%%%%%%%%
\subsection{Characteristic Values}
\label{subsec:characteristic-values-integral-equations}
%%%%%%%%%%%%%%%%%%%%%%%%%%%%%%%%%%%%%%%%%%%%%%%%%%%%%%%%%%%%

Depending on convention, the parameter \(\lambda\) for which

\[
u
=
\lambda Tu
\]

has a nonzero solution is called a characteristic value.

The corresponding nonzero \(u\) is an eigenfunction.

If

\[
Tu=\mu u,
\]

then

\[
\boxed{
\mu=\frac1\lambda.
}
\]

Care is needed because conventions differ across texts.

%%%%%%%%%%%%%%%%%%%%%%%%%%%%%%%%%%%%%%%%%%%%%%%%%%%%%%%%%%%%
\subsection{Symmetric Kernels}
\label{subsec:symmetric-integral-kernels}
%%%%%%%%%%%%%%%%%%%%%%%%%%%%%%%%%%%%%%%%%%%%%%%%%%%%%%%%%%%%

A real kernel is symmetric if

\[
\boxed{
K(x,t)
=
K(t,x).
}
\]

The associated integral operator behaves analogously to a symmetric
matrix under suitable hypotheses.

This leads to real eigenvalues and orthogonality properties.

%%%%%%%%%%%%%%%%%%%%%%%%%%%%%%%%%%%%%%%%%%%%%%%%%%%%%%%%%%%%
\subsection{Orthogonality of Eigenfunctions}
\label{subsec:integral-eigenfunction-orthogonality}
%%%%%%%%%%%%%%%%%%%%%%%%%%%%%%%%%%%%%%%%%%%%%%%%%%%%%%%%%%%%

For a sufficiently regular symmetric kernel, eigenfunctions corresponding
to distinct eigenvalues are orthogonal in the appropriate inner product.

For example,

\[
\boxed{
\int_a^b
u_m(x)u_n(x)\,dx
=
0,
\qquad
m\neq n.
}
\]

This parallels eigenvectors of symmetric matrices.

%%%%%%%%%%%%%%%%%%%%%%%%%%%%%%%%%%%%%%%%%%%%%%%%%%%%%%%%%%%%
\subsection{Hilbert--Schmidt Kernels Preview}
\label{subsec:hilbert-schmidt-kernels-preview}
%%%%%%%%%%%%%%%%%%%%%%%%%%%%%%%%%%%%%%%%%%%%%%%%%%%%%%%%%%%%

A kernel satisfying

\[
\boxed{
\int_a^b
\int_a^b
|K(x,t)|^2
\,dx\,dt
<
\infty
}
\]

defines an important class of compact integral operators on \(L^2\).

Such kernels are called Hilbert--Schmidt kernels.

Their spectral theory resembles finite-dimensional matrix theory in many
important ways.

%%%%%%%%%%%%%%%%%%%%%%%%%%%%%%%%%%%%%%%%%%%%%%%%%%%%%%%%%%%%
\subsection{Compact Operators Preview}
\label{subsec:compact-integral-operators-preview}
%%%%%%%%%%%%%%%%%%%%%%%%%%%%%%%%%%%%%%%%%%%%%%%%%%%%%%%%%%%%

Many integral operators with sufficiently regular kernels are compact.

Compact operators map bounded sets into sets whose closures are compact.

Their nonzero spectral values behave similarly to eigenvalues of finite
matrices:

\[
\boxed{
\text{nonzero eigenvalues are isolated},
}
\]

and can accumulate only at

\[
\boxed{
0
}
\]

under standard compact-operator theory.

%%%%%%%%%%%%%%%%%%%%%%%%%%%%%%%%%%%%%%%%%%%%%%%%%%%%%%%%%%%%
\subsection{Fredholm Alternative}
\label{subsec:fredholm-alternative}
%%%%%%%%%%%%%%%%%%%%%%%%%%%%%%%%%%%%%%%%%%%%%%%%%%%%%%%%%%%%

The Fredholm alternative describes solvability of equations such as

\[
\boxed{
(I-\lambda T)u=f.
}
\]

In a standard compact-operator setting, roughly one of the following
situations occurs:

\begin{enumerate}
    \item the corresponding homogeneous equation has only the trivial
          solution, and the nonhomogeneous equation has a unique
          solution for every suitable \(f\); or
    \item the homogeneous problem has nontrivial solutions, and the
          nonhomogeneous problem is solvable only when compatibility
          conditions are satisfied.
\end{enumerate}

%%%%%%%%%%%%%%%%%%%%%%%%%%%%%%%%%%%%%%%%%%%%%%%%%%%%%%%%%%%%
\subsection{Matrix Analogy for the Fredholm Alternative}
\label{subsec:fredholm-matrix-analogy}
%%%%%%%%%%%%%%%%%%%%%%%%%%%%%%%%%%%%%%%%%%%%%%%%%%%%%%%%%%%%

Compare

\[
\boxed{
A\mathbf{x}=\mathbf{b}.
}
\]

If \(A\) is invertible, there is a unique solution for every
\(\mathbf{b}\).

If \(A\) is singular, the homogeneous equation

\[
A\mathbf{x}=0
\]

has nonzero solutions, and the nonhomogeneous equation requires
compatibility with the left null space.

The Fredholm alternative generalizes this idea to operators.

%%%%%%%%%%%%%%%%%%%%%%%%%%%%%%%%%%%%%%%%%%%%%%%%%%%%%%%%%%%%
\subsection{Adjoint Kernel}
\label{subsec:adjoint-kernel-preview}
%%%%%%%%%%%%%%%%%%%%%%%%%%%%%%%%%%%%%%%%%%%%%%%%%%%%%%%%%%%%

For a complex-valued kernel, the formal adjoint kernel is

\[
\boxed{
K^*(x,t)
=
\overline{K(t,x)}.
}
\]

A kernel is self-adjoint when

\[
\boxed{
K(x,t)
=
\overline{K(t,x)}.
}
\]

For real kernels, this reduces to symmetry.

%%%%%%%%%%%%%%%%%%%%%%%%%%%%%%%%%%%%%%%%%%%%%%%%%%%%%%%%%%%%
\subsection{Fredholm Solvability Condition Preview}
\label{subsec:fredholm-solvability-condition-preview}
%%%%%%%%%%%%%%%%%%%%%%%%%%%%%%%%%%%%%%%%%%%%%%%%%%%%%%%%%%%%

When the homogeneous adjoint problem has nonzero solutions \(v\),
solvability of the nonhomogeneous equation can require

\[
\boxed{
\langle f,v\rangle=0.
}
\]

This is an infinite-dimensional analogue of orthogonality to the left
null space in linear algebra.

%%%%%%%%%%%%%%%%%%%%%%%%%%%%%%%%%%%%%%%%%%%%%%%%%%%%%%%%%%%%
\subsection{First-Kind Equations and Inverse Problems}
\label{subsec:first-kind-inverse-problems}
%%%%%%%%%%%%%%%%%%%%%%%%%%%%%%%%%%%%%%%%%%%%%%%%%%%%%%%%%%%%

A first-kind equation

\[
\boxed{
f(x)
=
\int_a^b
K(x,t)u(t)\,dt
}
\]

often asks us to recover \(u\) from smoothed observations \(f\).

Because integral operators can suppress high-frequency information,
inverting them may amplify noise.

Thus first-kind equations frequently define inverse problems.

%%%%%%%%%%%%%%%%%%%%%%%%%%%%%%%%%%%%%%%%%%%%%%%%%%%%%%%%%%%%
\subsection{Ill-Posedness}
\label{subsec:integral-equation-ill-posedness}
%%%%%%%%%%%%%%%%%%%%%%%%%%%%%%%%%%%%%%%%%%%%%%%%%%%%%%%%%%%%

An inverse problem is ill posed if one or more Hadamard properties fail:

\[
\boxed{
\text{existence},
}
\]

\[
\boxed{
\text{uniqueness},
}
\]

or

\[
\boxed{
\text{continuous dependence on the data}.
}
\]

For first-kind integral equations, instability is especially common.

%%%%%%%%%%%%%%%%%%%%%%%%%%%%%%%%%%%%%%%%%%%%%%%%%%%%%%%%%%%%
\subsection{Regularization Preview}
\label{subsec:integral-equation-regularization-preview}
%%%%%%%%%%%%%%%%%%%%%%%%%%%%%%%%%%%%%%%%%%%%%%%%%%%%%%%%%%%%

When direct inversion amplifies noise, one may replace the unstable
problem by a nearby stable optimization problem.

For example, Tikhonov regularization minimizes

\[
\boxed{
\|Tu-f\|^2
+
\alpha
\|u\|^2,
}
\]

where

\[
\alpha>0
\]

controls the amount of regularization.

This idea reappears in numerical analysis, inverse problems, and
statistical learning.

%%%%%%%%%%%%%%%%%%%%%%%%%%%%%%%%%%%%%%%%%%%%%%%%%%%%%%%%%%%%
\subsection{Abel Integral Equation Preview}
\label{subsec:abel-integral-equation-preview}
%%%%%%%%%%%%%%%%%%%%%%%%%%%%%%%%%%%%%%%%%%%%%%%%%%%%%%%%%%%%

An Abel-type equation has form

\[
\boxed{
f(x)
=
\int_0^x
\frac{
u(t)
}{
(x-t)^\alpha
}
\,dt,
\qquad
0<\alpha<1.
}
\]

The kernel has an integrable singularity at

\[
t=x.
\]

Such equations arise in inverse problems and connect with fractional
calculus.

%%%%%%%%%%%%%%%%%%%%%%%%%%%%%%%%%%%%%%%%%%%%%%%%%%%%%%%%%%%%
\subsection{Singular Integral Equations}
\label{subsec:singular-integral-equations}
%%%%%%%%%%%%%%%%%%%%%%%%%%%%%%%%%%%%%%%%%%%%%%%%%%%%%%%%%%%%

A singular integral equation contains a kernel with a nonintegrable or
principal-value singularity.

A classical form involves

\[
\boxed{
\operatorname{p.v.}
\int
\frac{
u(t)
}{
t-x
}
\,dt.
}
\]

These equations are important in complex analysis, fluid mechanics,
elasticity, and boundary integral methods.

%%%%%%%%%%%%%%%%%%%%%%%%%%%%%%%%%%%%%%%%%%%%%%%%%%%%%%%%%%%%
\subsection{Cauchy Principal Value}
\label{subsec:cauchy-principal-value}
%%%%%%%%%%%%%%%%%%%%%%%%%%%%%%%%%%%%%%%%%%%%%%%%%%%%%%%%%%%%

For an interior singularity at \(t=x\), the Cauchy principal value is
defined symmetrically by a limiting process such as

\[
\boxed{
\operatorname{p.v.}
\int_a^b
\frac{
f(t)
}{
t-x
}
\,dt
=
\lim_{\varepsilon\to0^+}
\left[
\int_a^{x-\varepsilon}
\frac{f(t)}{t-x}\,dt
+
\int_{x+\varepsilon}^{b}
\frac{f(t)}{t-x}\,dt
\right],
}
\]

when the limit exists.

%%%%%%%%%%%%%%%%%%%%%%%%%%%%%%%%%%%%%%%%%%%%%%%%%%%%%%%%%%%%
\subsection{Hammerstein Integral Equations}
\label{subsec:hammerstein-integral-equations}
%%%%%%%%%%%%%%%%%%%%%%%%%%%%%%%%%%%%%%%%%%%%%%%%%%%%%%%%%%%%

A Hammerstein equation has a common nonlinear form

\[
\boxed{
u(x)
=
f(x)
+
\lambda
\int_a^b
K(x,t)
F(t,u(t))
\,dt.
}
\]

The kernel acts linearly on a nonlinear function of the unknown.

Such equations occur in nonlinear boundary-value problems.

%%%%%%%%%%%%%%%%%%%%%%%%%%%%%%%%%%%%%%%%%%%%%%%%%%%%%%%%%%%%
\subsection{Urysohn Integral Equations}
\label{subsec:urysohn-integral-equations}
%%%%%%%%%%%%%%%%%%%%%%%%%%%%%%%%%%%%%%%%%%%%%%%%%%%%%%%%%%%%

A Urysohn equation has form

\[
\boxed{
u(x)
=
f(x)
+
\lambda
\int_a^b
K(x,t,u(t))
\,dt.
}
\]

The kernel itself depends nonlinearly on \(u(t)\).

This is more general than the Hammerstein form.

%%%%%%%%%%%%%%%%%%%%%%%%%%%%%%%%%%%%%%%%%%%%%%%%%%%%%%%%%%%%
\subsection{Fixed-Point Formulation}
\label{subsec:integral-equation-fixed-point}
%%%%%%%%%%%%%%%%%%%%%%%%%%%%%%%%%%%%%%%%%%%%%%%%%%%%%%%%%%%%

Many nonlinear integral equations can be written

\[
\boxed{
u=\mathcal{F}(u).
}
\]

Thus solving the equation means finding a fixed point of an operator.

This connects integral equations with fixed-point theorems such as:

\[
\boxed{
\text{Banach contraction theorem},
}
\]

and

\[
\boxed{
\text{Schauder fixed-point theorem}.
}
\]

%%%%%%%%%%%%%%%%%%%%%%%%%%%%%%%%%%%%%%%%%%%%%%%%%%%%%%%%%%%%
\subsection{Contraction Mapping Approach}
\label{subsec:integral-contraction-mapping}
%%%%%%%%%%%%%%%%%%%%%%%%%%%%%%%%%%%%%%%%%%%%%%%%%%%%%%%%%%%%

Suppose

\[
\mathcal{F}
\]

satisfies

\[
\boxed{
\|
\mathcal{F}(u)
-
\mathcal{F}(v)
\|
\leq
q
\|u-v\|,
\qquad
0<q<1.
}
\]

Then repeated iteration

\[
\boxed{
u_{n+1}
=
\mathcal{F}(u_n)
}
\]

converges to the unique fixed point under the assumptions of the Banach
fixed-point theorem.

%%%%%%%%%%%%%%%%%%%%%%%%%%%%%%%%%%%%%%%%%%%%%%%%%%%%%%%%%%%%
\subsection{Integral Equations with Memory}
\label{subsec:integral-equations-memory}
%%%%%%%%%%%%%%%%%%%%%%%%%%%%%%%%%%%%%%%%%%%%%%%%%%%%%%%%%%%%

Volterra equations naturally model memory:

\[
\boxed{
u(t)
=
f(t)
+
\int_0^t
K(t,s)u(s)\,ds.
}
\]

The current value depends on the entire history

\[
\boxed{
0\leq s\leq t.
}
\]

The kernel determines how strongly past states influence the present.

%%%%%%%%%%%%%%%%%%%%%%%%%%%%%%%%%%%%%%%%%%%%%%%%%%%%%%%%%%%%
\subsection{Convolution Memory Kernels}
\label{subsec:convolution-memory-kernels}
%%%%%%%%%%%%%%%%%%%%%%%%%%%%%%%%%%%%%%%%%%%%%%%%%%%%%%%%%%%%

If

\[
K(t,s)
=
k(t-s),
\]

then the effect depends only on the elapsed time between past and present.

The equation becomes

\[
\boxed{
u(t)
=
f(t)
+
\int_0^t
k(t-s)u(s)\,ds.
}
\]

This time-translation structure makes Laplace transforms especially
effective.

%%%%%%%%%%%%%%%%%%%%%%%%%%%%%%%%%%%%%%%%%%%%%%%%%%%%%%%%%%%%
\subsection{Renewal Equation Preview}
\label{subsec:renewal-equation-preview}
%%%%%%%%%%%%%%%%%%%%%%%%%%%%%%%%%%%%%%%%%%%%%%%%%%%%%%%%%%%%

A renewal-type equation may be written

\[
\boxed{
u(t)
=
f(t)
+
\int_0^t
u(t-s)\,dF(s),
}
\]

or in density form,

\[
\boxed{
u(t)
=
f(t)
+
\int_0^t
u(t-s)k(s)\,ds.
}
\]

Such equations arise in probability, reliability, branching processes,
and population models.

%%%%%%%%%%%%%%%%%%%%%%%%%%%%%%%%%%%%%%%%%%%%%%%%%%%%%%%%%%%%
\subsection{Numerical Solution of Integral Equations}
\label{subsec:numerical-integral-equations}
%%%%%%%%%%%%%%%%%%%%%%%%%%%%%%%%%%%%%%%%%%%%%%%%%%%%%%%%%%%%

When an exact solution is unavailable, an integral can be approximated by
a quadrature rule.

Suppose

\[
\int_a^b
g(t)\,dt
\approx
\sum_{j=1}^{m}
w_jg(t_j).
\]

Then

\[
\int_a^b
K(x,t)u(t)\,dt
\]

can be approximated by

\[
\boxed{
\sum_{j=1}^{m}
w_j
K(x,t_j)
u(t_j).
}
\]

%%%%%%%%%%%%%%%%%%%%%%%%%%%%%%%%%%%%%%%%%%%%%%%%%%%%%%%%%%%%
\subsection{Nystr\"om Method}
\label{subsec:nystrom-method}
%%%%%%%%%%%%%%%%%%%%%%%%%%%%%%%%%%%%%%%%%%%%%%%%%%%%%%%%%%%%

Consider

\[
u(x)
=
f(x)
+
\lambda
\int_a^b
K(x,t)u(t)\,dt.
\]

Replace the integral by quadrature:

\[
\boxed{
u(x)
\approx
f(x)
+
\lambda
\sum_{j=1}^{m}
w_j
K(x,t_j)
u(t_j).
}
\]

Evaluate at the quadrature points \(x=t_i\):

\[
\boxed{
u(t_i)
=
f(t_i)
+
\lambda
\sum_{j=1}^{m}
w_j
K(t_i,t_j)
u(t_j).
}
\]

This produces a finite linear system.

%%%%%%%%%%%%%%%%%%%%%%%%%%%%%%%%%%%%%%%%%%%%%%%%%%%%%%%%%%%%
\subsection{Matrix Form of the Nystr\"om Method}
\label{subsec:nystrom-matrix-form}
%%%%%%%%%%%%%%%%%%%%%%%%%%%%%%%%%%%%%%%%%%%%%%%%%%%%%%%%%%%%

Let

\[
u_i=u(t_i),
\]

and define

\[
A_{ij}
=
w_jK(t_i,t_j).
\]

Then

\[
\boxed{
\mathbf{u}
=
\mathbf{f}
+
\lambda A\mathbf{u}.
}
\]

Hence,

\[
\boxed{
(I-\lambda A)\mathbf{u}
=
\mathbf{f}.
}
\]

Thus numerical quadrature converts the integral equation into a matrix
equation.

%%%%%%%%%%%%%%%%%%%%%%%%%%%%%%%%%%%%%%%%%%%%%%%%%%%%%%%%%%%%
\subsection{Worked Example: Simple Quadrature Approximation}
\label{subsec:worked-integral-quadrature}
%%%%%%%%%%%%%%%%%%%%%%%%%%%%%%%%%%%%%%%%%%%%%%%%%%%%%%%%%%%%

\begin{workedexamplebox}{Midpoint Approximation}

Suppose

\[
u(x)
=
1
+
\int_0^1
xt\,u(t)\,dt.
\]

Using a one-point midpoint rule,

\[
t_1=\frac12,
\qquad
w_1=1.
\]

Then

\[
u(x)
\approx
1
+
x
\left(
\frac12
\right)
u\left(\frac12\right).
\]

Set

\[
x=\frac12.
\]

Then

\[
u\left(\frac12\right)
=
1
+
\frac14
u\left(\frac12\right).
\]

Thus,

\[
\frac34
u\left(\frac12\right)
=
1.
\]

Therefore,

\[
u\left(\frac12\right)
=
\frac43.
\]

Hence the approximation is

\begin{answerbox}
\[
\boxed{
u(x)
\approx
1+\frac23x.
}
\]
\end{answerbox}

\end{workedexamplebox}

%%%%%%%%%%%%%%%%%%%%%%%%%%%%%%%%%%%%%%%%%%%%%%%%%%%%%%%%%%%%
\subsection{Collocation Method Preview}
\label{subsec:collocation-integral-equations}
%%%%%%%%%%%%%%%%%%%%%%%%%%%%%%%%%%%%%%%%%%%%%%%%%%%%%%%%%%%%

Assume an approximate solution

\[
\boxed{
u_N(x)
=
\sum_{j=1}^{N}
c_j\phi_j(x).
}
\]

Insert this into the integral equation and define a residual

\[
\boxed{
R_N(x).
}
\]

Collocation chooses points

\[
x_1,\ldots,x_N
\]

and requires

\[
\boxed{
R_N(x_i)=0,
\qquad
i=1,\ldots,N.
}
\]

This yields algebraic equations for the coefficients \(c_j\).

%%%%%%%%%%%%%%%%%%%%%%%%%%%%%%%%%%%%%%%%%%%%%%%%%%%%%%%%%%%%
\subsection{Galerkin Method Preview}
\label{subsec:galerkin-integral-equations}
%%%%%%%%%%%%%%%%%%%%%%%%%%%%%%%%%%%%%%%%%%%%%%%%%%%%%%%%%%%%

The Galerkin method also assumes

\[
u_N(x)
=
\sum_{j=1}^{N}
c_j\phi_j(x).
\]

Instead of forcing the residual to vanish at points, require

\[
\boxed{
\langle
R_N,
\phi_i
\rangle
=
0,
\qquad
i=1,\ldots,N.
}
\]

Thus the residual is orthogonal to the chosen approximation space.

%%%%%%%%%%%%%%%%%%%%%%%%%%%%%%%%%%%%%%%%%%%%%%%%%%%%%%%%%%%%
\subsection{Projection Methods}
\label{subsec:projection-methods-integral-equations}
%%%%%%%%%%%%%%%%%%%%%%%%%%%%%%%%%%%%%%%%%%%%%%%%%%%%%%%%%%%%

Galerkin, least-squares, and related methods approximate an
infinite-dimensional problem by projecting it onto a finite-dimensional
subspace.

This parallels:

\[
\boxed{
\text{basis expansion in linear algebra},
}
\]

and

\[
\boxed{
\text{finite-element methods for PDEs}.
}
\]

%%%%%%%%%%%%%%%%%%%%%%%%%%%%%%%%%%%%%%%%%%%%%%%%%%%%%%%%%%%%
\subsection{Conditioning}
\label{subsec:integral-equation-conditioning}
%%%%%%%%%%%%%%%%%%%%%%%%%%%%%%%%%%%%%%%%%%%%%%%%%%%%%%%%%%%%

A numerical solution can be sensitive to perturbations in the data.

For first-kind equations, small changes in \(f\) can sometimes produce
large changes in the reconstructed \(u\).

This is a conditioning problem.

A numerical method cannot remove instability inherent in the underlying
mathematical problem without additional assumptions or regularization.

%%%%%%%%%%%%%%%%%%%%%%%%%%%%%%%%%%%%%%%%%%%%%%%%%%%%%%%%%%%%
\subsection{Integral Equation Workflow}
\label{subsec:integral-equation-workflow}
%%%%%%%%%%%%%%%%%%%%%%%%%%%%%%%%%%%%%%%%%%%%%%%%%%%%%%%%%%%%

A practical workflow is:

\begin{enumerate}
    \item identify the unknown function;
    \item identify the kernel;
    \item classify the equation as Fredholm or Volterra;
    \item determine whether it is first or second kind;
    \item determine whether it is linear or nonlinear;
    \item check for a separable or degenerate kernel;
    \item check whether differentiation converts the equation into an ODE;
    \item check whether the kernel has convolution form;
    \item use Laplace transforms for suitable Volterra convolution
          equations;
    \item investigate successive approximation;
    \item determine whether a Neumann series converges;
    \item investigate eigenvalue structure for homogeneous equations;
    \item use Green's functions when the equation comes from a
          boundary-value problem;
    \item use quadrature or projection methods when analytic methods fail;
    \item examine stability and conditioning of the problem.
\end{enumerate}

%%%%%%%%%%%%%%%%%%%%%%%%%%%%%%%%%%%%%%%%%%%%%%%%%%%%%%%%%%%%
\subsection{Choosing a Solution Method}
\label{subsec:choosing-integral-equation-method}
%%%%%%%%%%%%%%%%%%%%%%%%%%%%%%%%%%%%%%%%%%%%%%%%%%%%%%%%%%%%

If

\[
K(x,t)
=
\sum_j
g_j(x)h_j(t),
\]

consider:

\[
\boxed{
\text{finite-dimensional reduction}.
}
\]

If the equation is Volterra and differentiable, consider:

\[
\boxed{
\text{differentiation}.
}
\]

If the kernel is

\[
k(t-\tau),
\]

consider:

\[
\boxed{
\text{Laplace transforms}.
}
\]

If

\[
|\lambda|\|T\|<1,
\]

consider:

\[
\boxed{
\text{Neumann series}.
}
\]

For numerical approximation, consider:

\[
\boxed{
\text{Nystr\"om, collocation, or Galerkin methods}.
}
\]

%%%%%%%%%%%%%%%%%%%%%%%%%%%%%%%%%%%%%%%%%%%%%%%%%%%%%%%%%%%%
\subsection{Common Errors}
\label{subsec:integral-equation-common-errors}
%%%%%%%%%%%%%%%%%%%%%%%%%%%%%%%%%%%%%%%%%%%%%%%%%%%%%%%%%%%%

\subsubsection{Confusing Fredholm and Volterra Equations}

Fredholm equations use fixed limits.

Volterra equations have a variable limit such as

\[
\boxed{
\int_a^x.
}
\]

%%%%%%%%%%%%%%%%%%%%%%%%%%%%%%%%%%%%%%%%%%%%%%%%%%%%%%%%%%%%
\subsubsection{Confusing First and Second Kind}

First kind:

\[
\boxed{
f=Tu.
}
\]

Second kind:

\[
\boxed{
u=f+\lambda Tu.
}
\]

The placement of the unknown outside the integral matters.

%%%%%%%%%%%%%%%%%%%%%%%%%%%%%%%%%%%%%%%%%%%%%%%%%%%%%%%%%%%%
\subsubsection{Treating the Kernel as the Unknown}

In a standard integral equation, \(K(x,t)\) is given.

The unknown is \(u\).

Inverse problems can involve unknown kernels, but that is a different
problem.

%%%%%%%%%%%%%%%%%%%%%%%%%%%%%%%%%%%%%%%%%%%%%%%%%%%%%%%%%%%%
\subsubsection{Differentiating Without Adjusting the Variable Upper Limit}

For

\[
\int_a^x
K(x,t)u(t)\,dt,
\]

one must use the Leibniz rule.

Both the boundary term and the \(x\)-dependence of the kernel may
contribute.

%%%%%%%%%%%%%%%%%%%%%%%%%%%%%%%%%%%%%%%%%%%%%%%%%%%%%%%%%%%%
\subsubsection{Ignoring Conditions Lost During Differentiation}

Differentiating an integral equation can produce an ODE with additional
solutions.

The original integral equation must still be used to recover the correct
initial or boundary condition.

%%%%%%%%%%%%%%%%%%%%%%%%%%%%%%%%%%%%%%%%%%%%%%%%%%%%%%%%%%%%
\subsubsection{Assuming Every Neumann Series Converges}

The formal series

\[
I+\lambda T+\lambda^2T^2+\cdots
\]

requires convergence conditions.

It is not automatically valid for arbitrary \(\lambda\).

%%%%%%%%%%%%%%%%%%%%%%%%%%%%%%%%%%%%%%%%%%%%%%%%%%%%%%%%%%%%
\subsubsection{Confusing Operator Eigenvalues and the Parameter Convention}

If

\[
u=\lambda Tu,
\]

then

\[
Tu
=
\frac1\lambda u.
\]

Thus the operator eigenvalue is \(1/\lambda\), not necessarily
\(\lambda\).

%%%%%%%%%%%%%%%%%%%%%%%%%%%%%%%%%%%%%%%%%%%%%%%%%%%%%%%%%%%%
\subsubsection{Assuming Symmetry Without Checking the Kernel}

A symmetric kernel satisfies

\[
\boxed{
K(x,t)=K(t,x)
}
\]

for a real-valued problem.

This condition cannot be inferred merely from a visually similar formula.

%%%%%%%%%%%%%%%%%%%%%%%%%%%%%%%%%%%%%%%%%%%%%%%%%%%%%%%%%%%%
\subsubsection{Assuming First-Kind Equations Are Numerically Stable}

Integral operators often smooth functions.

Their inverse can strongly amplify measurement error.

Regularization may be necessary.

%%%%%%%%%%%%%%%%%%%%%%%%%%%%%%%%%%%%%%%%%%%%%%%%%%%%%%%%%%%%
\subsubsection{Applying Laplace Transforms to a Non-Convolution Kernel as if It Were Convolution}

The identity

\[
\mathcal{L}\{k*u\}
=
KU
\]

requires the convolution structure

\[
k(t-\tau).
\]

A general kernel

\[
K(t,\tau)
\]

does not transform into a simple product.

%%%%%%%%%%%%%%%%%%%%%%%%%%%%%%%%%%%%%%%%%%%%%%%%%%%%%%%%%%%%
\subsubsection{Ignoring Singular Kernels}

A kernel such as

\[
\frac1{t-x}
\]

may require principal-value interpretation.

Ordinary integration rules may not apply.

%%%%%%%%%%%%%%%%%%%%%%%%%%%%%%%%%%%%%%%%%%%%%%%%%%%%%%%%%%%%
\subsubsection{Assuming Numerical Discretization Guarantees Accuracy}

A fine quadrature grid does not automatically fix an ill-posed problem.

Discretization error and instability must be analyzed separately.

%%%%%%%%%%%%%%%%%%%%%%%%%%%%%%%%%%%%%%%%%%%%%%%%%%%%%%%%%%%%
\subsection{Applications}
\label{subsec:integral-equation-applications}
%%%%%%%%%%%%%%%%%%%%%%%%%%%%%%%%%%%%%%%%%%%%%%%%%%%%%%%%%%%%

\begin{applicationbox}

\textbf{Differential Equations.}
Initial-value and boundary-value problems can often be rewritten as
integral equations.

\medskip

\textbf{Potential Theory.}
Laplace and Poisson problems can be represented using Green's functions
and boundary integral equations.

\medskip

\textbf{Elasticity.}
Integral equations model stresses, displacements, cracks, and boundary
interactions.

\medskip

\textbf{Fluid Mechanics.}
Potential flow, vortex methods, and boundary formulations lead naturally
to integral equations.

\medskip

\textbf{Electromagnetism.}
Integral formulations describe fields generated by distributed charges
and currents.

\medskip

\textbf{Probability.}
Renewal equations and first-passage problems often produce Volterra
integral equations.

\medskip

\textbf{Population Biology.}
Age-structured and memory-dependent population models can be formulated
as integral equations.

\medskip

\textbf{Inverse Problems.}
Imaging, tomography, deconvolution, and parameter recovery frequently
lead to first-kind integral equations.

\medskip

\textbf{Control and Systems.}
Convolution equations represent input-output systems with memory.

\end{applicationbox}

%%%%%%%%%%%%%%%%%%%%%%%%%%%%%%%%%%%%%%%%%%%%%%%%%%%%%%%%%%%%
\subsection{Exercises}
\label{subsec:integral-equation-exercises}
%%%%%%%%%%%%%%%%%%%%%%%%%%%%%%%%%%%%%%%%%%%%%%%%%%%%%%%%%%%%

\begin{exercise}[1]
Define an integral equation.
\end{exercise}

\begin{exercise}[1]
Define the kernel of an integral equation.
\end{exercise}

\begin{exercise}[1]
Distinguish Fredholm and Volterra integral equations.
\end{exercise}

\begin{exercise}[1]
Distinguish equations of the first and second kind.
\end{exercise}

\begin{exercise}[1]
Define a separable kernel.
\end{exercise}

\begin{exercise}[1]
Define an integral operator.
\end{exercise}

\begin{exercise}[1]
Define a Neumann series.
\end{exercise}

\begin{exercise}[1]
Define an iterated kernel.
\end{exercise}

\begin{exercise}[1]
Define a resolvent kernel.
\end{exercise}

\begin{exercise}[1]
Define a symmetric kernel.
\end{exercise}

\begin{exercise}[2]
Classify

\[
u(x)
=
1
+
\int_0^1
xt\,u(t)\,dt
\]

by type and kind.
\end{exercise}

\begin{exercise}[2]
Classify

\[
u(x)
=
x
+
\int_0^x
u(t)\,dt.
\]
\end{exercise}

\begin{exercise}[2]
Classify

\[
f(x)
=
\int_0^1
K(x,t)u(t)\,dt.
\]
\end{exercise}

\begin{exercise}[2]
Determine whether

\[
K(x,t)=x^2t
\]

is separable.
\end{exercise}

\begin{exercise}[2]
Determine whether

\[
K(x,t)=x+t
\]

is a degenerate kernel and express it as a finite sum of separable terms.
\end{exercise}

\begin{exercise}[2]
Solve

\[
u(x)
=
2
+
\int_0^x
u(t)\,dt.
\]
\end{exercise}

\begin{exercise}[2]
Solve

\[
u(x)
=
x^2
+
\int_0^x
u(t)\,dt
\]

by differentiation.
\end{exercise}

\begin{exercise}[2]
Determine whether the kernel

\[
K(x,t)=e^{x+t}
\]

is symmetric.
\end{exercise}

\begin{exercise}[2]
Determine whether

\[
K(x,t)=x+t
\]

is symmetric.
\end{exercise}

\begin{exercise}[3]
Solve

\[
u(x)
=
1
+
x
\int_0^1
u(t)\,dt
\]

by reducing the equation to an algebraic equation.
\end{exercise}

\begin{exercise}[3]
Solve

\[
u(x)
=
x
+
x^2
\int_0^1
t\,u(t)\,dt.
\]
\end{exercise}

\begin{exercise}[3]
Reduce an equation with kernel

\[
K(x,t)
=
x+t
\]

to a finite linear system.
\end{exercise}

\begin{exercise}[3]
Derive the first three successive approximations for

\[
u(x)
=
1
+
\lambda
\int_0^x
u(t)\,dt.
\]
\end{exercise}

\begin{exercise}[3]
Show that the resulting series gives

\[
u(x)=e^{\lambda x}.
\]
\end{exercise}

\begin{exercise}[3]
Derive the operator Neumann series formally from

\[
(I-\lambda T)u=f.
\]
\end{exercise}

\begin{exercise}[3]
Compute the second iterated kernel for

\[
K(x,t)=xt
\]

on \([0,1]\).
\end{exercise}

\begin{exercise}[3]
Derive the Leibniz differentiation formula used for Volterra equations.
\end{exercise}

\begin{exercise}[3]
Convert

\[
y'=xy,
\qquad
y(0)=1
\]

into a Volterra integral equation.
\end{exercise}

\begin{exercise}[3]
Apply Picard iteration to the previous equation.
\end{exercise}

\begin{exercise}[3]
Use Laplace transforms to solve

\[
u(t)
=
1
+
\int_0^t
e^{-(t-\tau)}
u(\tau)\,d\tau.
\]
\end{exercise}

\begin{exercise}[3]
Derive the formula

\[
U(s)
=
\frac{
F(s)
}{
1-\lambda K(s)
}
\]

for a convolution integral equation.
\end{exercise}

\begin{exercise}[3]
Explain how Green's functions convert boundary-value problems into
integral equations.
\end{exercise}

\begin{exercise}[3]
Show that a homogeneous equation

\[
u=\lambda Tu
\]

is an eigenvalue problem for \(T\).
\end{exercise}

\begin{exercise}[3]
Explain the matrix analogy behind the Fredholm alternative.
\end{exercise}

\begin{exercise}[3]
Explain why first-kind equations commonly represent inverse problems.
\end{exercise}

\begin{exercise}[3]
Explain why regularization can stabilize an inverse integral equation.
\end{exercise}

\begin{exercise}[3]
Convert a quadrature approximation of a Fredholm equation into a matrix
system.
\end{exercise}

\begin{exercise}[3]
Apply a two-point quadrature rule to a simple second-kind Fredholm
equation.
\end{exercise}

\begin{exercise}[4]
Prove convergence of a Neumann series under

\[
|\lambda|\|T\|<1.
\]
\end{exercise}

\begin{exercise}[4]
Study resolvent kernels for Fredholm equations.
\end{exercise}

\begin{exercise}[4]
Study convergence of successive approximation for Volterra equations.
\end{exercise}

\begin{exercise}[4]
Prove existence and uniqueness for a Volterra equation using a contraction
mapping.
\end{exercise}

\begin{exercise}[4]
Study Hilbert--Schmidt integral operators.
\end{exercise}

\begin{exercise}[4]
Prove orthogonality of eigenfunctions for symmetric kernels.
\end{exercise}

\begin{exercise}[4]
Study compact self-adjoint integral operators.
\end{exercise}

\begin{exercise}[4]
Study the spectral theorem for compact self-adjoint operators.
\end{exercise}

\begin{exercise}[4]
Study the Fredholm alternative in Hilbert spaces.
\end{exercise}

\begin{exercise}[4]
Derive adjoint solvability conditions.
\end{exercise}

\begin{exercise}[4]
Study Abel integral equations and their inversion.
\end{exercise}

\begin{exercise}[4]
Investigate singular integral equations and principal values.
\end{exercise}

\begin{exercise}[4]
Study Hammerstein integral equations.
\end{exercise}

\begin{exercise}[4]
Study Urysohn integral equations.
\end{exercise}

\begin{exercise}[4]
Investigate Schauder fixed-point methods for nonlinear integral
equations.
\end{exercise}

\begin{exercise}[4]
Study renewal equations in probability.
\end{exercise}

\begin{exercise}[4]
Investigate Tikhonov regularization for first-kind equations.
\end{exercise}

\begin{exercise}[4]
Study singular-value decomposition of compact integral operators.
\end{exercise}

\begin{exercise}[4]
Derive convergence properties of the Nystr\"om method.
\end{exercise}

\begin{exercise}[4]
Study collocation methods for integral equations.
\end{exercise}

\begin{exercise}[4]
Study Galerkin projection methods.
\end{exercise}

\begin{exercise}[4]
Develop boundary-element methods from boundary integral equations.
\end{exercise}

%%%%%%%%%%%%%%%%%%%%%%%%%%%%%%%%%%%%%%%%%%%%%%%%%%%%%%%%%%%%
\subsection{Conceptual Summary}
\label{subsec:integral-equation-summary}
%%%%%%%%%%%%%%%%%%%%%%%%%%%%%%%%%%%%%%%%%%%%%%%%%%%%%%%%%%%%

A general second-kind Fredholm equation is

\[
\boxed{
u(x)
=
f(x)
+
\lambda
\int_a^b
K(x,t)u(t)\,dt.
}
\]

A Volterra equation has a variable upper limit:

\[
\boxed{
u(x)
=
f(x)
+
\lambda
\int_a^x
K(x,t)u(t)\,dt.
}
\]

In operator notation,

\[
\boxed{
(I-\lambda T)u=f.
}
\]

A formal inverse is represented by the Neumann series

\[
\boxed{
(I-\lambda T)^{-1}
=
I+\lambda T+\lambda^2T^2+\cdots.
}
\]

For convolution equations,

\[
\boxed{
u=f+\lambda k*u,
}
\]

the Laplace transform gives

\[
\boxed{
U(s)
=
\frac{
F(s)
}{
1-\lambda K(s)
}.
}
\]

A homogeneous equation

\[
\boxed{
u=\lambda Tu
}
\]

is an eigenvalue problem.

Green's functions convert differential equations into integral
representations.

Numerical quadrature converts integral equations into finite systems:

\[
\boxed{
(I-\lambda A)\mathbf{u}
=
\mathbf{f}.
}
\]

Integral equations therefore combine

\[
\boxed{
\text{integration}
+
\text{operator theory}
+
\text{differential equations}
+
\text{transforms}
+
\text{numerical analysis}.
}
\]

%%%%%%%%%%%%%%%%%%%%%%%%%%%%%%%%%%%%%%%%%%%%%%%%%%%%%%%%%%%%
\subsection{Section Connections}
\label{subsec:integral-equation-connections}
%%%%%%%%%%%%%%%%%%%%%%%%%%%%%%%%%%%%%%%%%%%%%%%%%%%%%%%%%%%%

\chapterconnection{
Integral Equations provide an alternative formulation of many differential
and boundary-value problems. Volterra equations connect directly with
initial-value problems and Picard iteration, while Fredholm equations
connect with boundary-value problems, Green's functions, eigenvalue
theory, and compact operators. Convolution kernels connect integral
equations with the Laplace and Fourier transform methods developed in
Section~\ref{sec:transform-methods}. First-kind equations introduce
important inverse and regularization problems. The next section, Delay
Differential Equations, extends ordinary differential equations by
allowing present rates of change to depend explicitly on past states,
introducing memory, delayed feedback, and new stability phenomena.
}

%%%%%%%%%%%%%%%%%%%%%%%%%%%%%%%%%%%%%%%%%%%%%%%%%%%%%%%%%%%%
% SECTION SOURCE TRACKING
%%%%%%%%%%%%%%%%%%%%%%%%%%%%%%%%%%%%%%%%%%%%%%%%%%%%%%%%%%%%

%------------------------------------------------------------
% 9.6 Integral Equations
%------------------------------------------------------------
%
% Source basis:
%   - Standard linear integral-equation theory.
%   - Standard Fredholm and Volterra methods.
%   - Standard operator, transform, and numerical formulations.
%
% Used for:
%   - integral equation definitions;
%   - kernels;
%   - Fredholm and Volterra equations;
%   - equations of the first, second, and third kind;
%   - homogeneous and nonhomogeneous equations;
%   - linear and nonlinear equations;
%   - integral operators;
%   - linear-algebra analogy;
%   - separable and degenerate kernels;
%   - successive approximation;
%   - Neumann series;
%   - iterated kernels;
%   - resolvent kernels;
%   - differentiation of Volterra equations;
%   - ODE integral formulations;
%   - Picard iteration;
%   - convolution equations;
%   - Laplace-transform solutions;
%   - Green's-function formulations;
%   - eigenvalue integral equations;
%   - characteristic values;
%   - symmetric kernels;
%   - Hilbert--Schmidt kernels;
%   - compact operators;
%   - Fredholm alternative;
%   - adjoint kernels;
%   - first-kind inverse problems;
%   - ill-posedness and regularization;
%   - Abel and singular equations;
%   - Hammerstein and Urysohn equations;
%   - fixed-point formulations;
%   - memory kernels and renewal equations;
%   - quadrature methods;
%   - Nystrom method;
%   - collocation;
%   - Galerkin methods;
%   - numerical conditioning;
%   - applications and exercises.
%
% Notes:
%   - Delay Differential Equations is developed next in Section 9.7.
%   - Differential-Algebraic Equations completes Chapter 9 in
%     Section 9.8.
%   - Numerical integral-equation methods are developed more deeply in
%     Chapter 11.
%
%%%%%%%%%%%%%%%%%%%%%%%%%%%%%%%%%%%%%%%%%%%%%%%%%%%%%%%%%%%%